\section{High dimensional isotropic Gaussian random function simulation}

If one wants to perform Gradient descent on a high dimensional random
function, then it is necessary to simulate the gradients together
with the function values. This is possible, since the covariance
of the gradient to the function value is
\[
	\Cov(\nabla \rf(x), \rf(y)) = \E[\nabla \rf(x) \rf(y)]
	= \nabla_x \E[\rf(x)\rf(y)] = \nabla_x \C(x,y),
\]
where we assumed zero mean \(\E[\rf(x)]=0\) for simplicity.
And similarly
\[
\Cov(\nabla \rf(x), \nabla\rf(y)) = \nabla_x \nabla_y\C(x,y).
\]
So the covariance of the derivatives is there derivative
of the covariance function. This makes simulation
of gradients possible, however a gradient contains \(\dims\)
entries if \(x\in \real^\dims\). Each of these is a random variable.
So the vector \(\nabla \rf(x_1), \dots, \nabla\rf(x_n)\) contains
\(n\dims\) random variables, its covariance matrix is of size \((n\dims)^2\)
and its cholesky decomposition costs \(\bigO((n\dims)^3)\). This is
prohibitively expensive for \(\dims \gg 1\).

For isotropic Gaussian random functions this computational complexity may
however be reduced to 
\[
	\bigO((n \min\{n, \dims\})^3),
\]
using the ``Adatped-Span'' from
\citep{benningGradientSpanAlgorithms2024}. In large dimension the above is of
course \(\bigO(n^6)\). Here we will give a more detailed explanation.

An ``isotropic'' Gaussian random function (GRF) \(\rf\) has mean and covariance functions
of the form
\[
	\mu(x) = \mu(\tfrac{\norm{x}^2}2)
	\quad\text{and}\quad \C(x,y) = \kappa(\tfrac{\norm{x}^2}2, \tfrac{\norm{y}^2}2, \scp{x, y})
\]
Of course if something is a function of \(\tfrac{\norm{x}^2}2\), it can also be written
as a function of \(\norm{x}\). However this definition interacts better with
taking derivatives, which is necessary as outlined above.

A special case are the ``stationary isotropic`` GRF that has mean and covariance
of the form (for a list of examples see \citep[Ch.~5]{rasmussenGaussianProcessesMachine2006})
\[
	\mu(x) = \mu \in \real
	\quad \text{and} \quad \C(x,y) = C(\tfrac{\norm{x-y}^2}2).
\]
In the following we will assume isotropy!

Observe that because an orthonormal basis change does not change the covariance
we essentially have that the partial derivatives all have the same distribution
-- well almost: The partial derivatives of \(\rf(x)\) in directions that are
orthogonal ot \(x\) all have the same distribution. Moreover, if they are
orthogonal they are uncorrelated and since we are in the Gaussian case thereby
independent!

In general, if
\[
	V_n = \Span\set{x_0, \dots, x_n, \nabla \rf(x_0),\dots, \nabla \rf(x_{n-1})}
\]
then \textbf{any orthogonal basis} \(w_{\dim(V_n)}, \dots, w_{\dims-1}\) of the
orthogonal complement \(V_n^\perp\)  of \(V_n\) has the property that
\(D_{w_i} \rf(x_n)\) are independent identically distributed!

So how does the covariance look in the first step? Well \(V_0 = \set{x_0}\)
and taking any orthonormal basis \(w_i\) of \(V_0^\perp\) we have
\[
	\Cov\Bigl[\begin{pmatrix}
		\rf(x_0)
		\\
		D_{x_0}\rf(x_0)
		\\
		D_{w_1} \rf(x_0)
		\\
		\vdots
		\\
		D_{w_{\dims-1}}\rf(x_0)
	\end{pmatrix}
	\Bigr]
	= \begin{pmatrix}
		* & *\\
		* & *\\
		& & \sigma^2 \\
		& & & \ddots
		\\
		& & & &\sigma^2 
	\end{pmatrix},
\]
where \(\sigma^2\) is the variance of each \(D_{w_i}\rf(x_0)\). The key to
understand now is that a basis change of the \(w_i\) does not change this
matrix. We can thus sort of choose a ``future basis'' of \(V_\dims\) before
knowing it, because we do not know yet what points \(x_i\) the user will choose.

This is what the \texttt{LazyAdaptedSpan} does. It represents \(V_n\)
and its future extension. And it also tries to avoid ever
realizing this basis. By expressing everything in terms of the coordinates of
\(V_n\) it is thereby possible to avoid a lot of the computational cost.
This is what the code \verb|x0 = grf.into_adapted_span(x0)| does. It turns a
numpy vector into a coordinate vector with respect to the Adaptive span of
the GRF. Note that optimizers such as ADAM have coordinate specific learning
rates which depends on the standard basis. They can therefore not work
with the adapted span. In this case the adapted span needs to be realized which
costs \(\bigO(n^2\dims)\) due to the Gram-Schmidt orthogonalization under the hood.
In that case the total complexity is
\[
\bigO(n^6 + n^2\dims)
\]
This should still be reasonable for moderate dimension sizes.

Observe that as \(V_n\) grows with \(n\), the non-diagonal part grows with every
step. This means that the covariance matrices will essentially look like this:
\begin{center}
	\includegraphics*[width=\textwidth]{images/matrix_layout.jpg}
\end{center}
Let \(v_0, \dots, v_{\dims-1}\) be an orthonormal basis of \(V_n\) and
for \(n\) large enough such that \(\dim V_n = \dims\). Then for
\[
	\Sigma_{kl} \coloneq \Cov\Bigl[\begin{pmatrix}
		\rf(x_k)	
		\\
		D_{v_0} \rf(x_k)
		\\
		\vdots
		\\
		D_{v_{\dims-1}} \rf(x_k)
	\end{pmatrix}, \begin{pmatrix}
		\rf(x_l)	
		\\
		D_{v_0} \rf(x_l)
		\\
		\vdots
		\\
		D_{v_{\dims-1}} \rf(x_l)
	\end{pmatrix}
	\Bigr]
\]
we have that the covariance matrix of all the values so far is given by
\[
	\begin{bmatrix}
		\Sigma_{11} & \cdots & \Sigma_{1n}
		\\
		\vdots & & \vdots
		\\
		\Sigma_{n1} & \cdots & \Sigma_{nn},
	\end{bmatrix}
\]
where each \(\Sigma_{lk}\) is a ``Kite-matrix'' with dense upper left
corner of size
\[
	\dim V_{\max\set{l, k}},
\]
and a diagonal in the bottom right, where every entry is the same (a ``scaled identity'').

Considering the first section on the simulation of Gaussian random vectors it is
therefore vital that we can perform the Cholesky decomposition
of Block matrices, where each block is a ``Kite-matrix''. And the block matrix
must be stored in such a way as to make appending rows easy. For the
covariance matrix we can use symmetry and store
\[
	[\Sigma_{11}, (\Sigma_{21}, \Sigma_{22}), (\Sigma_{31}, \dots, \Sigma_{33}), \dots),]
\]
for the Cholesky decomposition we can also only store the bottom left corner as
the upper right is zero anyway.
\textcolor{red}{
	However if we want to mix and match the decompositions (SVD, eigh, cholesky) as outlined in the section
	before, this becomes more delicate. Maybe it is necessary to add another
	block structure in between, which may contain multiple Kite matrices.
}

\section{TODO}

\begin{enumerate}
	\item Add more covariance examples. You can see above why it is necessary to
	be able to take gradients.
	\item First add code to profile performance,
	\item Then try to write better performing (non-python) matrix computations underneath.
	\item Also testing: It is really difficult to write random simulation code,
	because the outcome is by definition random. So you cannot state exactly what you
	expect. But you can plot the random function from the simulation and
	compare it to previous implementation to see if they pass the smell test.
	The matrix decompositions are also deterministic and it is sometimes possible
	to compare them directly to a less sophisticated reference implementation
	that is simpler to understand but less performant.
\end{enumerate}

